% =============================================================================
% thesis_template.tex
% Workshop: Scientific Writing with LaTeX, Overleaf & Prism
% Purpose:  A thesis / dissertation template showing front matter, chapters,
%           appendices, and bibliography.  Compiles as a standalone file.
%           In a real project you would split chapters into separate files
%           and use \input{chapters/ch1.tex}, etc.
% =============================================================================

\documentclass[12pt, a4paper, oneside]{book}
% "book" class provides \chapter, \frontmatter, \mainmatter, \appendix, etc.

% --- Packages -----------------------------------------------------------------
\usepackage[utf8]{inputenc}
\usepackage[T1]{fontenc}
\usepackage{amsmath, amssymb, amsthm}
\usepackage{graphicx}
\usepackage{booktabs}
\usepackage{hyperref}
\usepackage{geometry}
\geometry{margin=2.5cm}
\usepackage{setspace}          % Line spacing control
\onehalfspacing                % 1.5 line spacing (common thesis requirement)
\usepackage{fancyhdr}          % Custom headers and footers
\usepackage{lipsum}            % Dummy text for demonstration — remove in real use

% --- Header / footer style ----------------------------------------------------
\pagestyle{fancy}
\fancyhf{}
\fancyhead[L]{\leftmark}
\fancyhead[R]{\thepage}
\renewcommand{\headrulewidth}{0.4pt}

% --- Theorem environments -----------------------------------------------------
\theoremstyle{plain}
\newtheorem{theorem}{Theorem}[chapter]
\newtheorem{lemma}[theorem]{Lemma}
\theoremstyle{definition}
\newtheorem{definition}[theorem]{Definition}

% ==============================================================================
\begin{document}

% ##############################################################################
%  FRONT MATTER
% ##############################################################################
\frontmatter   % Pages numbered with roman numerals; chapters are un-numbered

% --- Title Page ---------------------------------------------------------------
\begin{titlepage}
    \centering
    \vspace*{2cm}
    {\LARGE\bfseries Title of the Thesis\par}
    \vspace{1.5cm}
    {\Large Subtitle or Expanded Title (if any)\par}
    \vspace{2cm}
    {\large
        A thesis submitted in fulfilment of the requirements\\
        for the degree of\\[0.5em]
        \textbf{Doctor of Philosophy}\\[0.5em]
        in the\\[0.5em]
        Department of Mathematics\\
        Faculty of Science\\
        University Name
    \par}
    \vfill
    {\large
        \textbf{Author Name}\\[0.4em]
        Supervisor: Prof.\ Supervisor Name\\[1em]
        \the\year
    \par}
\end{titlepage}

% --- Declaration (optional) ---------------------------------------------------
\chapter*{Declaration}
\addcontentsline{toc}{chapter}{Declaration}

I declare that this thesis is my own unaided work.  It is being submitted for
the degree of Doctor of Philosophy at University Name.  It has not been
submitted before for any degree or examination at any other university.

\vspace{2cm}
\noindent
\rule{6cm}{0.4pt}\\
Author Name\\
\today

% --- Abstract -----------------------------------------------------------------
\chapter*{Abstract}
\addcontentsline{toc}{chapter}{Abstract}

\lipsum[1]   % Replace with your actual abstract

% --- Acknowledgements ---------------------------------------------------------
\chapter*{Acknowledgements}
\addcontentsline{toc}{chapter}{Acknowledgements}

I would like to thank my supervisor, Prof.\ Supervisor Name, for their
guidance and support.  I am also grateful to my colleagues at the Department
of Mathematics for many fruitful discussions.

\lipsum[2]   % Replace with your own text

% --- Table of Contents, Lists ------------------------------------------------
\tableofcontents
\listoffigures
\listoftables

% ##############################################################################
%  MAIN MATTER
% ##############################################################################
\mainmatter    % Arabic page numbers restart from 1; chapters are numbered

% In a real thesis, replace the content below with:
%   \input{chapters/introduction}
%   \input{chapters/literature_review}
%   ...

% --- Chapter 1 ----------------------------------------------------------------
\chapter{Introduction}
\label{ch:introduction}

% \input{chapters/introduction}   % Uncomment when using separate files

\section{Background}
\label{sec:background}
\lipsum[3-4]

\section{Research Questions}
\label{sec:research_questions}
The main questions addressed in this thesis are:
\begin{enumerate}
    \item What are the convergence properties of method X?
    \item How does method X compare to existing approaches?
    \item Can the results be extended to infinite-dimensional settings?
\end{enumerate}

\section{Outline of the Thesis}
\label{sec:outline}
This thesis is organised as follows.
Chapter~\ref{ch:literature} reviews the relevant literature.
Chapter~\ref{ch:main_results} presents the main results.
Chapter~\ref{ch:conclusion} concludes with a summary and future work.

% --- Chapter 2 ----------------------------------------------------------------
\chapter{Literature Review}
\label{ch:literature}

% \input{chapters/literature_review}

\section{Classical Results}
\label{sec:classical}
\lipsum[5]

\begin{definition}[Contraction Mapping]
    \label{def:contraction}
    A mapping $T \colon X \to X$ on a metric space $(X, d)$ is a
    \emph{contraction} if there exists $0 \leq q < 1$ such that
    $d(Tx, Ty) \leq q \, d(x, y)$ for all $x, y \in X$.
\end{definition}

\begin{theorem}[Banach Fixed-Point Theorem]
    \label{thm:banach}
    Every contraction on a complete metric space has a unique fixed point.
\end{theorem}

\section{Recent Developments}
\label{sec:recent}
\lipsum[6-7]

% --- Chapter 3 ----------------------------------------------------------------
\chapter{Main Results}
\label{ch:main_results}

% \input{chapters/main_results}

\section{Problem Formulation}
\label{sec:formulation}
\lipsum[8]

\begin{lemma}
    \label{lem:key}
    Under assumptions A1--A3, the operator $T$ satisfies
    \begin{equation}
        \lVert T^{n} x - x^{*} \rVert \leq q^{n} \lVert x - x^{*} \rVert.
        \label{eq:contraction_bound}
    \end{equation}
\end{lemma}

\begin{proof}
    The result follows by induction on $n$ using
    Definition~\ref{def:contraction}.
\end{proof}

\section{Numerical Experiments}
\label{sec:numerical}
\lipsum[9]

% Example figure placeholder
\begin{figure}[htbp]
    \centering
    % \includegraphics[width=0.7\textwidth]{figures/convergence_plot}
    \rule{0.7\textwidth}{5cm}   % Placeholder — remove when using real image
    \caption{Convergence of the iterative method for test problem 1.}
    \label{fig:convergence}
\end{figure}

% Example table
\begin{table}[htbp]
    \centering
    \caption{Comparison of iteration counts for different methods.}
    \label{tab:comparison}
    \begin{tabular}{lccc}
        \toprule
        \textbf{Method} & \textbf{Problem 1} & \textbf{Problem 2}
                         & \textbf{Problem 3} \\
        \midrule
        Newton           &  6  &  8  &  5 \\
        Modified Newton  & 12  & 15  & 10 \\
        Proposed         &  5  &  7  &  4 \\
        \bottomrule
    \end{tabular}
\end{table}

% --- Chapter 4 ----------------------------------------------------------------
\chapter{Conclusion and Future Work}
\label{ch:conclusion}

% \input{chapters/conclusion}

\section{Summary of Contributions}
\lipsum[10]

\section{Future Research Directions}
\begin{itemize}
    \item Extend the convergence analysis to Hilbert space settings.
    \item Investigate adaptive step-size strategies.
    \item Apply the method to large-scale PDE-constrained optimisation.
\end{itemize}

% ##############################################################################
%  APPENDICES
% ##############################################################################
\appendix      % Subsequent \chapter commands produce Appendix A, B, ...

\chapter{Proof of Auxiliary Lemma}
\label{app:aux_proof}

% \input{chapters/appendix_a}

\lipsum[11]

\begin{align}
    \lVert x_{n+1} - x^{*} \rVert
        &\leq \lVert T(x_n) - T(x^{*}) \rVert \notag \\
        &\leq q \lVert x_n - x^{*} \rVert.
\end{align}

\chapter{Supplementary Data}
\label{app:data}

% \input{chapters/appendix_b}

Additional tables and figures that support the results in
Chapter~\ref{ch:main_results} are collected here.

\lipsum[12]

% ##############################################################################
%  BACK MATTER
% ##############################################################################
\backmatter    % Un-numbered chapters again (for bibliography, index, etc.)

% --- Bibliography -------------------------------------------------------------
% For a real thesis, use BibTeX or BibLaTeX:
%   \bibliographystyle{plain}
%   \bibliography{references}

\begin{thebibliography}{99}

\bibitem{banach1922}
    S.~Banach,
    \emph{Sur les op\'{e}rations dans les ensembles abstraits et leur
    application aux \'{e}quations int\'{e}grales},
    Fundamenta Mathematicae, \textbf{3} (1922), 133--181.

\bibitem{kantorovich1948}
    L.~V.~Kantorovich,
    \emph{On Newton's method for functional equations},
    Doklady Akad.\ Nauk SSSR, \textbf{59} (1948), 1237--1240.

\bibitem{ortega1970}
    J.~M.~Ortega and W.~C.~Rheinboldt,
    \emph{Iterative Solution of Nonlinear Equations in Several Variables},
    Academic Press, New York, 1970.

\end{thebibliography}

\end{document}
